\section{Metaphysical Positivism}

\begin{quotex}
We, on the contrary---basing ourselves on a tradition much more ancient and real than the one which can be claimed by the ``faith'' of Western man, on a tradition which is not proved by doctrines, but by deeds and acts of power and clairvoyance---affirm instead the possibility and the concrete reality of what we have called ``Wisdom''. We thus assert the possibility of a positive, direct, methodical, empirical knowledge in the ``metaphysical'' field, just as science strives to gain in the physical field, and, just like science, it remains above any moral or philosophical belief of men.

\flright{\textsc{Julius Evola}, \emph{Those Who Know and Those Who Believe}\footnote{\url{https://www.gornahoor.net/PaganImperialism/ThoseWhoKnow.htm}} in \textit{Pagan Imperialism}\footnote{\url{https://www.gornahoor.net/PaganImperialism/PaganImperialism.htm}}}
\end{quotex}


In the aftermath of the destruction of the Traditional world order resulting from the French Revolution, the philosopher \textbf{Auguste Comte} sought to re-establish it on a positive basis. Thus, Positivism would accept as ``real'' only what can be experienced and the laws that are discovered in that experience. It would supersede the theological and ideological speculations of earlier stages of humanity\footnote{\url{https://gornahoor.net/?p=530\%20}}.\footnote{We use the word ``ideological'' where Comte uses ``metaphysical'', since it more accurately reflects his intention. We therefore reserve ``metaphysics'' as the science of the suprasensible.}

Now, Evola emphasizes that what he means by ``metaphysics'' is just as exact, methodical and based on the experience of the real, as any empirical science. And just like Positivism, it is independent of any theological or ideological speculations. However, unlike Positivism, it is not the result of a progression or evolution of humanity, but rather a recovery of the most ancient or \textbf{Primordial Tradition}. Yet this is not the accumulation over the centuries of a set of doctrines that can be learned and memorized, but rather something that must be recovered anew by each man through his own efforts and vision. As such it is neither transcendent faith nor physical science, but

\begin{quotex}
necessarily a surpassing of both religious unrealism and materialised realism by a transcendent, virile, Olympian positivism.
\end{quotex}


Evola describes this more specifically:

\begin{quotex}
we maintain the possibility of forms of experience different from the sensory forms of the common man, not ``given'', not ``normal'', which can be reached by means of certain active processes of inner transformation. The peculiarity of such transcendent experiences ... is to be direct, concrete, and individual, as much as sensory experience itself, and yet to see reality, beyond the contingent, spatio-temporal aspect characteristic of everything that is sensory.

\flright{\textsc{Julius Evola}, \emph{Science vs Wisdom}\footnote{\url{https://www.gornahoor.net/PaganImperialism/ScienceVsWisdom.htm}} in \textit{Pagan Imperialism}\footnote{\url{https://www.gornahoor.net/PaganImperialism/PaganImperialism.htm}}}
\end{quotex}


As we pointed out, this involves transcending consensus reality, as described by Josephin Peladan\footnote{\url{https://gornahoor.net/?p=995}}. Evola puts it this way:

\begin{quotex}
the whole so-called ``problem of knowledge'' is enclosed within the interiority of every being, and does not depend on ``culture'', but on his capacity for freeing himself from the human, i.e., from the sensory, the rational, and the emotional, and of identifying himself with one or another form of ``metaphysical'' experience.
\end{quotex}


This is the critical, yet difficult, point to accept: \emph{viz.}, the metaphysical principle that ``to know is to be''. Unless and until a man has regenerated his consciousness to the Primordial State, he does not and cannot know anything about it. Unless and until he becomes a True Man\footnote{\url{https://gornahoor.net/?p=1039}}, he can only speculate about it.

\begin{quotex}
To know, according to Wisdom, does not mean ``to think'', but to be the thing known: to live it, to realise it inwardly. One does not really know a thing unless one can actively transform one's consciousness into it. Therefore, only what ensues from direct individual experience will count as knowledge. And, this is just the opposite of the modern mentality, for which, whatever appears immediately to the individual is called ``phenomenon'', or ``subjective'', and so it posits some other thing behind it as ``true reality'', which is simply imagined or presumed (the ``thing in itself'' of the philosophers, the ``Absolute'' of vulgar religion, ``matter'', ``ether'', or ``energy'' of science). Wisdom is an absolute positivism which regards only what can be grasped by direct experience as real, and everything else as unreal, abstract, and illusory.


\flright{\textsc{Julius Evola}, \emph{Science vs Wisdom}\footnote{\url{https://www.gornahoor.net/PaganImperialism/ScienceVsWisdom.htm}} in \textit{Pagan Imperialism}\footnote{\url{https://www.gornahoor.net/PaganImperialism/PaganImperialism.htm}}}
\end{quotex}



\flrightit{Posted on 2010-11-17 by Cologero}

\begin{center}* * *\end{center}

\begin{footnotesize}
\begin{sffamily}
\texttt{KarlH on 2010-11-17 at 23:45 said:}

This is perhaps the most frustrating and reward paradox of our pursuit. The person who achieves a Primordial State can know things, but not necessarily in a communicable sense. They can leave hints and cause people to consider a transformation of self, but it's nearly impossible to speak of the inner experience of transformation. This does, admittedly, read like lunacy in the soulless positivism of our modern age.

\hfill

\texttt{Sati Shankar on 2010-11-18 at 09:15 said:}

In that state of awareness, it is true that we only have unspeakable
... but before that every effort is mental and there lies the

territory of metaphysical positivism... though, seemingly an attempt to unify the dual... .brave an attempt... .all the best.

\hfill

\texttt{James O'Meara on 2010-11-18 at 10:13 said:}

This, or rather similar passages on Experience in the book Magic, are what I consider to be the `essential Evola' in the sense that everything else depends on this. Evola ``sees'' the positivist demand for experience and ``raises'' the stakes by observing that the kind of experience one has depends on ones state of consciousness, or being; higher, or lower. Rather than seeking to transcend experience by something below, or behind, it, or fleeing into mathematical abstractions, Evola proposes staying with experience and changing ones consciousness. The positivist may call this lunacy, but if he refuses to follow the methods offered and ``see'' for himself, what of it?

\hfill

\texttt{James O'Meara on 2010-11-18 at 10:19 said:}

Coincidentally, just stumbled across this on Dmitri Orlov's post-economic collapse site. For `incompetent', read: positivist.

``A key insight is offered by the Dunning-Krueger effect, defined and experimentally tested by Justin Kruger and David Dunning at Cornell University. Kruger and Dunning proposed that, ``for a given skill, incompetent people will:

tend to overestimate their own level of skill;
fail to recognize genuine skill in others;
fail to recognize the extremity of their inadequacy;
recognize and acknowledge their own previous lack of skill, if they can be trained to substantially improve.''

``Krueger and Dunning, and other experimenters, have shown this effect to be quite pronounced. Competent people initially assumed that others were competent as well, and were able to correct their misperception once they were allowed to examine the work of others. Incompetent people, on the other hand, were only able to recognize competence in others after being taught to recognize their own incompetence. Thus, a weaker version of point 4 above suffices: incompetent people do not need to become competent, but to be able to judge the superior competence of others they do have to gain some insight into their own incompetence.''

\hfill

\texttt{Cologero on 2010-11-18 at 18:26 said:}

It's a dicey problem, whose solution requires some openness and perhaps a little humility. Rudolf Steiner makes an interesting point when he complains that those who deny the infallibility of the Pope, make themselves into little popes themselves. A common criticism made of him is this: ``only that can be true which a simple, unpretentious mind can perceive and understand.'' To which Steiner responds:

\begin{quotex}
which means, of course, only what the speaker can perceive and understand. He therefore requires that no one should perceive and understand anything different from what he himself perceives and understands. The infallibility of the Pope is not acknowledged in such assemblies, the infallibility of the individual is claimed instead.
\end{quotex}


Unfortunately, neither Steiner in his organisation, nor Evola in the Magic groups, were able to develop a consistent training method that would lead to a common body of reproducible, esoteric knowledge. This is where it differs from the positive sciences.

\hfill

\texttt{James O'Meara on 2010-11-19 at 11:09 said:}

I would think the ``common body of reproducible, esoteric knowledge'' would be the body of Tradition, or the traditions, in their theoretical and practical wings. It's right there to be explored and used, and conclusions drawn for oneself, if one wants, particularly today, which may not be altogether a good thing in the big picture.

Of course, attempts like Evola's have the disadvantage of being undertaken in singularly poor circumstances, in both time and space.

The Steiner quote is ironic, since while of course rejecting the ``acceptable to the simple mind'' approach, which would hardly allow anything beyond basic arithmetic, to say nothing of quantum mechanics or modern medicine, he never the less seemed to think he had developed a relatively simple technique that `anyone' could learn, and more importantly use to gain knowledge, without any ``traditional data'' as Guenon would say. This is his ``modern man has grown up and needs no helping hand'' mentality, like all modern cults inspired by ``modern science.'' Thus he and his disciples were never able to distinguish genuine clairvoyance perhaps from personal fantasies. Wilson or Lachman give as an example his detailed accounts of Arthur at Tintagel which later studies show to have been impossible --- no castle in Arthur's time . I suppose this explains Evola's rather equivocal attitude to Steiner, while the more theoretical Guenon condemns him outright.

Michael Hoffman at egodeath.com who seems to have disappeared was much interested in the lousy track record of ``meditation based'' approaches to enlightenment, especially such as modern Boomer Buddhism. Humility indeed: how many have ever been enlightened? He favored intense drug-induced crisis, interspersed with study under the guidance of the already experienced; think: Eleusis.

\hfill

\texttt{James O'Meara on 2010-11-19 at 11:23 said:}

More to the point: the teachings are known, the methods are there, like modern science. Like modern science, they seem to be too hard for the `average Joe' to just grasp my reading a few books. Unlike modern science, which can turn out thousands of doctors and engineers, the `success rate' is close to zero. Some say the methods are sound, but modern man is not fit Guenon or simply needs different methods Evola in his Tantric modes , others question the validity of the methods themselves Hoffman, promoting drugs over meditation .

\hfill

\texttt{Will on 2010-11-20 at 11:30 said:}

Good comments, gentlemen.

There is also the problem of recognizing the genuine article when it's there. There are all sorts of preconceived ideas about what an enlightened being will be like. For example, most Boomer Buddhists will dismiss the idea that Evola had any degree of realization, because if he had he would have become a liberal humanist, right? I recall talking to an aged hippie who had spent a lot of time in various ``spiritual'' groups, and he said that he only ever met one really enlightened teacher. His criteria? The guy ran a soup kitchen for the homeless, whereas none of the others had their own 501c charity.

Of course, those on the right have their own preconceived ideas. Presumably the only ones who don't are those who are themselves enlightened.

Perhaps the idea of enlightened vs. unenlightened (and hence, arguing over who is which) should be abandoned in favor of a more concrete approach of recognizing superior qualities. Too many people want be on the top of the hierarchy, but it seems to me that we do better to try to make an honest assessment of where we're at in it, and try to move up from wherever that may be. If you want to learn how to play guitar and you're a complete beginner, it doesn't really matter if your teacher is a master or a novice; either can teach you the basics, which is where you have to start.

To continue the analogy, mastering a musical instrument takes decades of consistent practice, and most highly accomplished musicians will tell you that even at their high level of attainment, they are still learning and growing. I think spiritual practice is like this as well. Even if there is an end point where one can say, like the Buddha, ``Done, what needed to be done,'' it will be slow progress up to that point. Promises of quick enlightenment, whether through drugs or ``secret methods'' or whatever, are mostly sales pitches. But the traditions tease us by saying that there are indeed a small number of people who can attain profound realization very quickly because of their superior aptitude. And who doesn't like to entertain the possibility that he is one of them?

\hfill

\texttt{James O'Meara on 2010-11-20 at 11:49 said:}

Will, I was just about to post this and saw your comment. Perhaps this is a more ``concrete'' way to look at it, as opposed to `enlightened' and `unenlightened', or as previously suggested, a more `humble' approach.

Evola in Doctrine of Awakening talks about `determining the vocations.' Only the Aryan Man in the Buddhist sense, not Himmler's will respond to the realization of samsara with a noble disdain rather than a desperate, purely reactionary asceticism and follow the Noble Path. Others will respond by conceiving of some guilt and seeking forgiveness through submission to the divine will e.g., Christianity or saying ``Well, let's eat and drink and make merry'' the Epicurean, or the modern Dawkinsian, the materialist in general. And so on.

The Superior Man is not interested in `arguing' or `proving' his point of view the `demon of dialectics' as discussed in a previous chapter . Various men simply respond according to type, and there is nothing more to say. The Superior Man simply draws his conclusions, and seeks to work on himself.

\hfill

\texttt{James O'Meara on 2010-11-20 at 11:56 said:}

Will,

Soup kitchen spirituality, indeed. Perhaps related to the idea, often announced by smug `secularists,' that ``of course, Jesus was a wise and progressive rabbi of his time.'' I notice they love to throw in the `rabbi' part, an odd combination of `see, I'm not an anti-Semite' with belittling just some guy hired by the congregation, not the Son of God at the same time.

I used to be amused or disgusted at East West Books here in NYC now closed, I guess a victim of the recession which was a huge store stuffed with every imaginable crystal, yoga mat, aryavedic nostrum and hippie geegaw, and thousands of books of various levels of seriousness, and off in one corner, one book by Evola; the Metaphysics of Sex, of course, chosen no doubt for its title if they had read it, they would have chucked it out in the trash ; as for Guenon, etc., not a trace. Complete works of Annie Besant, though!

\hfill

\texttt{Will on 2010-11-21 at 11:06 said:}

James,

I've seen Evola's books in a couple new age bookstores as well. The stores must have standing orders with Inner Traditions and their other imprints, and the books just get sent to them. I don't know how else to explain it. I can't imagine the book ordering guy perusing the table of contents of Revolt Against the Modern World or Men Among the Ruins and saying, ``Yes, this will go next to Eckhart Tolle and Deepak Chopra.''

In regards to Jesus as a progressive rabbi, you might be interested in a pastor named Mark Driscoll from the Pacific Northwest, whose ministry is very reactionary against that sort of thing, not so much from an anti-Jewish perspective as anti-feminist. He's somehow managed to make Calvinism cool to Seattle hipsters, which might qualify as a miracle in and of itself. Part of the appeal is that his approach is unabashedly masculine. I wouldn't consider him a representative of Tradition, but I do think he's interesting as a cultural and religious phenomenon.

\url{http://www.youtube.com/watch?v=lex6orNNzTs}


\hfill

\texttt{James O'Meara on 2010-11-22 at 09:52 said:}

Will,

You're probably right about Evola on IT standing orders. They probably get one copy which never sells and just stays on the shelf.

I will have to look into this Pastor. Calvin for hipsters? I'm reminded of Death to the World, a group of `punk monks' bringing Orthodoxy to `a dying Protestant scene'' (\url{http://www.deathtotheworld.com/about/about.html}). It's also interesting to see the masculine angle. Evola actually was able to say a few good words for `sola fide' as being a way to bring home the seriousness of what's at stake, forcing an existential choice, that was relatively valid for a certain human type.

\end{sffamily}
\end{footnotesize}

